% Answers to Chapter 5's exercises, lectures/L05/appendix/b_exercises.md. The book's own answers:
% the lecture publishes none.
\answerchapter{5}{Timers}

\answer{c5:ex:1}
The prescaler divides the CPU clock, and the tick period is the inverse of the tick frequency:
\[
  f_\text{tick} = \frac{16\,000\,000\,\text{Hz}}{256} = 62\,500\,\text{Hz},
  \qquad
  T_\text{tick} = \frac{1}{62\,500\,\text{Hz}} = 16\,\text{\textmu s} = 0.016\,\text{ms},
\]
the 16\,\textmu s of the \textdiv256 row in the table of \secref{c5:a:2}.

\answer{c5:ex:2}
An 8-bit timer overflows every 256 ticks:
\[
  T_\text{ovf} = 256 \times 16\,\text{\textmu s} = 4.096\,\text{ms},
  \qquad
  N = \frac{250\,\text{ms}}{4.096\,\text{ms}} = 61.04.
\]
Rounded to the nearest whole number the way the chapter's \code{OVF_MAX} does it,
$\lfloor 61.04 + 0.5 \rfloor = 61$ overflows. Check: $61 \times 4.096\,\text{ms} =
249.856$\,ms, 0.144\,ms (0.06\,\%) short of 250\,ms. (Counted in ticks it is the same:
$250 / 0.016 = 15\,625$ ticks, and $15\,625 / 256 = 61.04$.)

\answer{c5:ex:3}
Normal mode with overflow counting can only make whole multiples of the overflow period, 256 ticks
for an 8-bit timer (16.384\,ms at \textdiv1024), and it takes an interrupt for every overflow even
though only the last one does anything. Any other period is rounded to the nearest multiple, an
error of up to half an overflow period. CTC chooses the period in single ticks: the counter is
cleared when it reaches \code{OCRnA}, so the period is $\text{\code{OCRnA}} + 1$ ticks, with
one interrupt per period and no counting in software. Reach for CTC when the period has to be
accurate, is not a multiple of the overflow period, or is short enough that the interrupt overhead
matters; overflow counting is fine when roughly right is right enough.

Yes: 100\,ms is one. At \textdiv1024, Timer 0 overflows every 16.384\,ms, and
$100 / 16.384 = 6.10$, so the best Normal mode can do is 6 overflows, 98.304\,ms, 1.7\,\% short.
Timer 1 in CTC with \code{OCR1A = 1562} gives $1563 \times 64\,\text{\textmu s} = 100.032$\,ms,
0.03\,\% off. A period that is a whole number of ticks is hit exactly: the 2\,s of
Exercise~\ref{c5:ex:6} is 31,250 ticks.

\answer{c5:ex:4}
Only \code{TIMEOUT_MS} changes, from \code{500U} to \code{250U}: 2\,Hz is one on/off cycle every
500\,ms, so the LED must toggle every 250\,ms. The prescaler, \code{TICK_PERIOD_MS},
\code{TICK_MAX}, \code{OVF_PERIOD_MS} and the ISR all stay as they are, and \code{OVF_MAX} follows
from the macro:
\[
  \text{\code{OVF_MAX}} = \left\lfloor \frac{250}{0.064 \times 256} + 0.5 \right\rfloor
    = \lfloor 15.26 + 0.5 \rfloor = 15.
\]
Check: $15 \times 16.384\,\text{ms} = 245.76$\,ms per toggle, a 491.52\,ms period, so the LED
blinks at 2.03\,Hz. (16 overflows would give 262.144\,ms, further from 250.)

\answer{c5:ex:5}
\textbf{Inside the ISR.} \code{ovf_count} is a \code{static} local of the ISR, so no other code can
even name it, and on the AVR the ISR runs to completion with interrupts disabled, so nothing can
change it between two of the ISR's own accesses. The compiler sees every access in one function
body. It still has to load the value when the ISR starts and store it before the ISR returns,
because a \code{static} outlives the call, and between those two points whatever it keeps in a
register is correct. \code{volatile} only forbids optimizations that are already safe here, so
removing it changes nothing about what the program does.

\textbf{Waited on from \code{main()}.} Now the loop \code{while (10U > ovf_count) {}} reads the
variable, and its body neither writes it nor calls anything. Without \code{volatile}, the compiler
may assume nothing else writes it either: it has no idea that an ISR exists. It is allowed to read
\code{ovf_count} once, before the loop, and test that register copy forever; if the first read is
below 10, \code{main()} hangs in an endless loop while the ISR keeps incrementing the variable in
memory. Whether that happens depends on the build, not on the code: at \code{-O0} every access is
compiled as a load from memory and the loop works, while with optimization, as at the \code{-Os}
the book builds with, avr-gcc hoists the load out of the loop and it hangs. A different compiler
version, a change in inlining, or a function call added to the loop body can each make the bug
appear or disappear again. \code{volatile} makes every test of \code{ovf_count} in the loop a fresh
read from memory.

\answer{c5:ex:6}
The number of ticks in 2\,s, less one, because the counter takes $\text{\code{OCR1A}} + 1$
ticks to reach the compare value:
\[
  \text{ticks} = \frac{2000\,\text{ms}}{0.064\,\text{ms}} = 31\,250,
  \qquad
  \text{\code{OCR1A}} = 31\,250 - 1 = 31\,249.
\]
It fits: Timer 1 is 16-bit, so \code{OCR1A} holds up to 65,535, while an 8-bit timer's 255 could
not. Check: $(31\,249 + 1) \times 64\,\text{\textmu s} = 2.000$\,s exactly, since 16\,MHz divided by
1024 is exactly 15,625 ticks per second. The chapter's macros agree, and show why they round: in
\code{float} arithmetic \code{2000U / 0.064F - 1U} comes to 31248.998, which a plain cast would
truncate to 31248; adding 0.5 first gives 31249.
